% Prima pagina del CV: profilo discorsivo che apre il documento.
% La includono gli scheletri in latex/tex/ (per ora solo cv_full.tex).
\vspace{15pt}
\begin{wrapfigure}{l}{6.5cm} % 'l' = left, 'r'
  \vspace{-12pt}               % move image up
  \includegraphics[width=6cm,height=6.78cm]{profile_2.png} % square image
    \vspace{-16pt}                % adjust bottom spacing
\end{wrapfigure}

Enrico Mensa is currently enrolled as a research fellow (RTDa) at the Computer Science Department of the University of Turin, where he earned his Master’s Degree in Computer Science (2016: summa cum laude, with mention for the outstanding CV) and his PhD, also in Computer Science (2020: passed with honor). 

\paragraph{Research.} EM is a member of the bigger Artificial Intelligence group of the Department, namely the Content-Centered Computing research group. His research is mostly in Natural Language Processing (NLP), knowledge representation, lexical resources and natural language semantics, although EM has set research collaborations in further areas, such as in event extraction, human-computer interaction, cultural heritage, as witnessed by his publications.
EM is actively involved in projects focused on the analysis of natural language (mostly from the cultural heritage area, the legal domain, and the medical field), and has started collaborations with research groups both at the national and international level. 
EM authored over 35 refereed works, published in conference/workshop proceedings and journals, has served as a reviewer some of the most prominent journals and conferences in the NLP landscape. He also recently started his activity as an editor with a special issue for the Frontiers in Artificial Intelligence Series ('Information Extraction for Health Documents').

\paragraph{Teaching.} EM teaches the course `Introduzione all'Informatica e alle Reti' at the Department of Cultures, Politics and Society. He also has been providing lectures for masters (in Artificial Intelligence, and Data Science and Business Intelligence), and in has previously served as a tutor for courses in the Bachelor degree in Computer Science  (First and second course in Programming, and Development of Software Applications). EM supervised nine research theses (four master's and five bachelor's thesis) all on topics from his research lines. He is also a PhD advisor for one PhD student.

\paragraph{Scientific Dissemination.} EM is active in AI scientific dissemination, with invitations to speak at prominent national events such as the Rumore Festival (Fanpage.it). He curates a YouTube channel on artificial intelligence that has garnered over 53,000 subscribers and 3,000,000 views through educational content on various AI topics.

\paragraph{Professional duties and service.} EM served as a representative of PhD students within the Departmental Governance, and is now active as a representative of research fellows in the Department board.

%Representative of the PhD students (2017-2019??) and representative of the research fellows in the governance of the CS Department.

\vspace{30pt}